\chapter{Introduction}
\label{ch:introduction}

\section{Motivation}
\label{sec:motivation}

German tax advisors (\textit{Steuerberater}), auditors
(\textit{Wirtschaftspr\"ufer}), and lawyers (\textit{Rechtsanw\"alte}) process
large volumes of client documents --- incoming invoices
(\textit{Eingangsrechnungen}) and supporting receipts (\textit{Belege}) --- as a
daily core activity. Recent progress in document \acp{VLM} makes it tempting to
automate this work with cloud services. For these professions, however, the
question is not only one of data protection: it also touches professional
secrecy, which is protected by criminal law.

% SUPERVISOR GATE (issue #96): this paragraph was deliberately SOFTENED from an
% earlier, overstated version that claimed cloud processing is simply forbidden.
% The nuance below is verified against the federal chamber of tax advisors'
% guidance and German legal/IT commentary, but a legal claim in a thesis needs
% human sign-off. Confirm the wording with the supervisor before submission.
The \ac{DSGVO} can be satisfied for cloud processing through a data-processing
agreement and EU data residency. Professional secrecy adds a second, separate
requirement: under \S203~\ac{StGB} together with \S62a~\ac{StBerG}, any external
service provider who gains access to client data must be bound to
confidentiality in text form, and foreign providers must additionally satisfy a
comparable-secrecy test \cite{bstbk2026}. This requirement is not categorically
unmeetable: at least one major cloud provider offers a contractual amendment
specifically for holders of professional secrets. Two considerations nevertheless
keep local processing attractive. First, the providers of the strongest
general-purpose language models do not currently offer an equivalent amendment.
Second, and independently of any contract, providers under United States
jurisdiction remain subject to extraterritorial access powers, which is precisely
the residual risk the comparable-secrecy test addresses. Local inference removes
the question rather than managing it: the document never leaves the firm. This
thesis takes that constraint as its operational premise --- not as its research
contribution --- and asks what is technically achievable \emph{within} it.

\section{Problem Statement}
\label{sec:problem}

Working locally confines the system to models that fit on commodity, on-premises
hardware. This raises a concrete empirical question: how reliably do small,
locally hosted, open-weights document \acp{VLM} --- the wave of specialised models
released from late~2025 onward --- convert German B2B invoices into a structured,
legally meaningful form? Prior work on invoice processing with vision models
\cite{berghaus2025} benchmarked large cloud models but did not include the recent
specialised open-weights document \acp{VLM}; that omission is the gap this thesis
addresses.

A second problem emerges only once such a system is measured seriously. An
end-to-end extraction score conflates three very different failure sources: the
model may fail to \emph{read} the page, it may read correctly but fail to
\emph{structure} what it read, or the \emph{measurement itself} may be unfair ---
penalising an answer that is in fact correct as printed. Separating these is a
prerequisite for spending effort in the right place, and it turns out to change the
engineering conclusion substantially.

\section{Research Questions}
\label{sec:research-questions}

The guiding question of this thesis is:

\begin{quote}
\itshape
How reliably can current open-weights, locally hosted document \acp{VLM} convert
German business-to-business invoices into a legally meaningful structured
representation on commodity on-premises hardware --- and which stage of that
pipeline limits the achievable reliability?
\end{quote}

\noindent It is addressed through four sub-questions:

\begin{enumerate}
  \item How do recent open-weights document \acp{VLM} compare with one another on
        German e-invoice field extraction?
  \item Does splitting the task into a reading stage and a structuring stage
        outperform a single model that goes from page image straight to
        structured output?
  \item What accuracy is attainable inside the memory and throughput envelope of
        commodity local hardware?
  \item Where does the residual error actually originate --- in reading, in
        structuring, or in the measurement?
\end{enumerate}

\noindent These are made falsifiable through a hypothesis set registered before
the corresponding data was collected; Chapter~\ref{ch:methodology} states the
hypotheses and the pre-registration discipline. Part of that set concerns system
layers that were subsequently placed outside the scope of this thesis; those
hypotheses are reported as registered-but-not-evaluated in
Chapter~\ref{ch:limitations-future-work} rather than quietly dropped.

\section{Contributions}
\label{sec:contributions}

% TODO(contributions): finalise once results land. Draft claims:
This thesis contributes:

\begin{itemize}
  \item an evaluation of recent open-weights local document \acp{VLM} on German
        e-invoice extraction, closing the cloud-only gap in \cite{berghaus2025};
  \item a compliance-aware evaluation suite that weights fields by their legal
        importance under \S14~\ac{UStG}/\S33~\ac{UStDV} rather than treating all
        fields alike under a single token \ac{F1};
  \item a \emph{measurement-validity} methodology: a documented chain of
        ground-truth and scoring corrections that separates genuine model error
        from scoring artefact --- for instance, cases where the model correctly
        reproduces what the page shows while the reference stores a different
        but equivalent representation;
  \item an \emph{attribution} method that decomposes end-to-end error into a
        reading share, a structuring share and a measurement share, using
        reference-quality transcripts as an upper bound on what the structuring
        stage can achieve; and
  \item a locally deployable prototype and a reproducible evaluation harness as
        engineering artefacts.
\end{itemize}

\section{Scope and Delimitations}
\label{sec:scope}

This thesis is \emph{not} a startup pitch, a pure \ac{OCR} benchmark, or a
legal-tech treatise: professional secrecy frames the motivation but is not the
contribution. It does not claim to beat commercial systems in general; it
characterises a German-specific cell of the problem space where the local-only
constraint binds.

The evaluated scope is the document-reading and structuring layer. A knowledge
graph over the extracted data and natural-language querying over that graph were
designed (Chapter~\ref{ch:system-design}) but deliberately not built or evaluated;
so was a direct comparison against cloud models. Chapter~\ref{ch:limitations-future-work}
states what was left out and why.

\paragraph{A note on terminology.} \emph{\ac{OCR}-free} in this thesis is an
architectural statement, not a claim about model names: the pipeline contains no
separate character-recognition engine, and the vision-language model maps page
pixels directly to output tokens. Several models in the evaluated cohort
nevertheless carry \ac{OCR} in their names, because their authors named them
after the task they perform rather than after the architecture they use. The two
usages are independent, and the naming is not a contradiction.

\section{Thesis Structure}
\label{sec:structure}

% Keep in sync with main.tex's include list.
Chapter~\ref{ch:background} introduces the technical and legal background, and
Chapter~\ref{ch:related-work} positions the thesis against prior work.
Chapter~\ref{ch:system-design} presents the system design and states which parts
of it are evaluated here. Chapter~\ref{ch:methodology} specifies the corpus, the
ground truth, the metric suite and the pre-registration discipline.
Chapter~\ref{ch:results} reports the results, and
Chapter~\ref{ch:implementation} describes the prototype.
Chapter~\ref{ch:discussion} interprets the findings,
Chapter~\ref{ch:limitations-future-work} sets out the limitations and the work
left undone, and Chapter~\ref{ch:conclusion} concludes.
